\section{Results}
\label{sec:results}

\subsection{How Large Is the Small-County Gap?}
\label{sec:gap}

Figure~\ref{fig:gradient}a plots raw participation by population decile. The share of county--declarations with any HMGP subapplication rises from 14.3 percent in the smallest decile (median population 3,426) to 43.2 percent in the largest (median 413,325), and the share with a funded project follows almost the same path. The two lines are close because most decided subapplications are approved, so the differences between counties emerge before a request reaches the state or FEMA. The rise is steepest in the upper deciles, but participation also increases, with small reversals, across the lower half of the distribution.

Table~\ref{tab:baseline} reports Equation~(\ref{eq:baseline}) with controls added across columns. Comparing counties within the same state and year (column 1), small counties are 11.4 percentage points less likely to obtain any funded HMGP project (Panel B). Replacing state-by-year effects with declaration fixed effects (column 2) barely changes the estimate, and adding the hazard and damage controls (column 3) reduces it only to 10.5 points. The gap is therefore not explained by small counties experiencing weaker storms. Adding income, poverty, race, and nonmetropolitan status reduces the gap to 8.8 percentage points (column 4), or 32 percent of the larger-county mean of 27.1 percent. The gap in filing any subapplication (Panel A) is of the same size, 8.9 percentage points on a mean of 26.1 percent, an early sign that the shortfall occurs before any decision on a request is made.

The gap is proportionally larger for investments aimed at tornado risk. Small counties are 3.1 percentage points less likely to obtain a safe room or wind retrofit under the declaration (Panel C, column 4), against a larger-county mean of 6.3 percent, so their probability is about half that of larger counties. Over a three-year window that includes projects funded under later declarations, the gap is 8.1 percentage points on a mean of 14.6 percent (Table~\ref{tab:robust}).

The gap is a gradient rather than a threshold (Figure~\ref{fig:gradient}b). Relative to the largest population decile, and within declarations with full controls, the probability of a funded project is 35.5 percentage points lower in the smallest decile, 23.1 points lower in the fifth, and 11.5 points lower in the ninth (Appendix Table~\ref{tab:deciles}); each decile lies below the one above it. In continuous form, a one-log-point increase in population, roughly a 2.7-fold difference, is associated with a 7.8 percentage point higher probability of a funded project (Table~\ref{tab:robust}). Differences in participation thus arise across the whole range of county sizes and are not confined to a few very small counties.

\subsection{How Much Reflects the Number of Potential Applicants?}
\label{sec:structure}

A county's claim on an HMGP allocation is made by the entities located in it, and small counties have fewer of them. Column (5) of Table~\ref{tab:baseline} adds the log numbers of municipalities and townships, special districts, and school districts. The gap in funded projects falls from 8.8 to 5.0 percentage points, and the gap in subapplications from 8.9 to 5.4 points. About two-fifths of the small-county gap is thus statistically accounted for by the number of entities that could apply. The count of school districts has the largest coefficient of the three (Appendix Table~\ref{tab:controls}), consistent with schools serving as sites for community safe rooms. For safe rooms and wind retrofits, local-government counts account for a smaller share of the gap (from 3.1 to 2.4 percentage points), leaving a gap equal to 38 percent of the larger-county mean. Both residual gaps are precisely estimated, with standard errors of at most one percentage point.

This component reflects the institutional structure of local government in small counties rather than the behavior of any single government. Because the counts are not assigned independently of county size, the decomposition describes how the gap is distributed across this channel; it does not show what would happen if a county gained or lost local governments.

\subsection{Does the Gap Arise at Application or at Approval?}
\label{sec:pipeline}

Table~\ref{tab:pipeline} and Figure~\ref{fig:pipeline} follow requests through the grant process, holding constant the column (5) controls, including local-government counts. Small counties are 5.4 percentage points less likely to file any subapplication and file about 4.7 percent fewer. Among county--declarations with at least one decided subapplication (2,515 observations), the approval rate of small counties differs from that of larger counties by $-0.5$ percentage points (95 percent confidence interval from $-4.2$ to $3.2$) on a mean of 86.6 percent, and the withdrawal rate by $-0.9$ points ($-4.4$ to $2.7$). The denial rate is 1.4 percentage points higher ($p = 0.03$), from a low base of 1.7 percent; this difference is statistically detectable but too small to account for a meaningful part of the gap. At the subapplication level, small-county requests are smaller and report lower benefit--cost ratios, but neither difference is statistically distinguishable from zero once local-government counts are controlled (Table~\ref{tab:pipeline}, note). The evidence supports H2: the shortfall occurs at the decision to apply, not in the review of submitted requests. This comparison of decision rates relies mainly on recent records, because few older closed or denied subapplications carry the county information needed for linkage (Section~\ref{sec:limitations}).

The county government provides the cleanest check on the role of institutional structure, because every county has exactly one. With the full set of controls, small county governments are 2.4 percentage points less likely to file an HMGP subapplication under a declaration, on a mean of 11.1 percent (column 3), and 2.4 points less likely to receive a project, on a mean of 11.4 percent (column 7). A gap of about one-fifth for a single, comparable applicant cannot be attributed to the number of entities, which supports H3.

\subsection{Why Do Small Counties Apply Less?}
\label{sec:explanations}

Table~\ref{tab:mechanisms} assembles four tests of explanations for the application gap.

\textit{Response to hazard events.} Panel A estimates Equation~(\ref{eq:shock}). In larger counties, a declaration in which a tornado crossed a populated place raises the number of subapplications by 4.3 percent ($p = 0.02$) and the probability that the county government applies by 2.4 percentage points ($p = 0.045$), relative to other declarations designating the same county. In small counties the corresponding responses are $-2.9$ percent and $-1.8$ points, and the differences between small and larger counties are $-7.2$ percent ($p = 0.046$) and $-4.2$ points ($p = 0.06$). Figure~\ref{fig:shock} displays these responses. The continuous measure, log residents in tornado paths, gives the same pattern ($0.35$ points per log unit for larger county governments, $p = 0.02$; interaction $-0.60$, $p = 0.05$). For the broader outcomes, any subapplication and any funded project, the interactions have the same sign but are imprecise ($p = 0.23$ and $p = 0.16$). Because the county fixed effects remove persistent differences between places, this pattern is difficult to reconcile with the view that small counties simply need less mitigation: larger counties respond to a damaging event with more applications, while small counties show no comparable response. With $p$-values near conventional thresholds and several outcomes examined, this evidence for H4 is suggestive rather than conclusive.

\textit{Fiscal and administrative resources.} Panel B considers fiscal and administrative capacity measured per resident. Small county governments raise more tax revenue per resident and spend more per resident on financial administration and central staff than larger ones, consistent with the fixed cost of operating a government being spread over fewer residents. Controlling for these measures leaves the gap in subapplications at 5.4 percentage points, and the gap is not larger where per-resident administrative spending is low (interaction $-0.5$ points, standard error 1.6). Holding population constant, absolute administrative spending and tax revenue also do not predict applications. The ability to fund the 25 percent match, at least as captured by these fiscal measures, does not explain the gap. These measures record spending, not the specialized expertise that applications require, so this result does not rule out a role for capacity more broadly.

\textit{State delivery.} Panel C asks whether states that file applications on behalf of localities have smaller gaps. Using the share of each state's HMGP subapplications filed by state agencies before 2009 (mean 18.5 percent), a one-standard-deviation higher state share is associated with a 2.6 percentage point smaller gap in subapplications ($p = 0.04$) and a 3.0 point smaller gap in funded projects ($p = 0.02$) in 2009--2021. A comparable measure for regional planning and development bodies, which file few subapplications in any state, shows no association. The state-share result is a cross-state association, and states that file more applications may differ in other ways, but it is consistent with survey evidence that state assistance shapes local access \parencite{smith2020state}.

\textit{Reimbursement of management costs.} Panel D and Figure~\ref{fig:drra} test whether reimbursing subrecipients' management costs narrowed the gap. In raw terms, the gap in subapplications was 16.2 percentage points for declarations made from 2010 to July 2017 and 12.4 points afterward. Within counties and declarations, however, the change is $-1.9$ percentage points (standard error 3.6) for subapplications and $-1.7$ points (3.6) for funded projects, and the placebo comparison at August 2014 is likewise null. The event-study coefficients fluctuate around zero before and after the reform without a break (smallest $p = 0.14$). The estimates are imprecise: the upper end of the 95 percent confidence interval for subapplications corresponds to a narrowing of about 5 percentage points, so a modest effect cannot be excluded, but a narrowing of the size of the full baseline gap is not supported. A smaller exploratory comparison around FEMA's pre-calculated benefits for residential safe rooms, whose start date we cannot establish precisely, also shows no differential change for small counties (Appendix~S3).

\subsection{Robustness and Qualifications}
\label{sec:robustness}

Table~\ref{tab:robust} reports the gap in funded projects under alternative choices. Clustering by state increases the standard error only modestly for the main outcome. Restricting the sample to county--declarations in which a tornado crossed a populated place, where demand for wind mitigation should be highest, enlarges the gap to 14.3 percentage points on a mean of 33.8 percent and to 9.2 points for safe rooms and retrofits on a mean of 14.5 percent. Limiting the sample to declarations that included Individual Assistance leaves the gap in any funded project unchanged ($-8.8$ points) and somewhat enlarges the gap in safe rooms and retrofits ($-4.3$ points). Dropping each state in turn yields gaps between 8.4 and 9.4 percentage points for any funded project and between 2.7 and 3.5 points for safe rooms and retrofits. The gap over a five-year window is 10.5 points on a mean of 19.7 percent.

Two checks qualify the results. When statewide HMGP projects under the same declaration are credited to every designated county, the gap in any project nearly disappears ($-0.9$ points). This outcome is close to saturated, however, because about 89 percent of larger counties' county--declarations, and 86 percent of small counties', are covered by some statewide project, mostly planning and similar activities. For safe rooms and retrofits the gap remains ($-2.8$ points on a mean of 9.9 percent), so state-run projects offset only part of the shortfall in local investment. Among the 361 county--declarations under declarations whose incident type is recorded as a tornado rather than a severe storm, the estimates are small and imprecise ($-1.5$ points, standard error 4.0); this subsample is too small to yield informative estimates.
